\section{Project Overview}
In this project, we develop a teleoperation system for the UR10e 6-DOF manipulator using an iPhone as the input device. The phone's pose, tracked by Apple's ARKit framework, drives the robot through two control modes:
\begin{itemize}
    \item \textbf{Velocity mode}: The phone's instantaneous twist (linear and angular velocity) is mapped to joint velocities via resolved-rate control using the geometric Jacobian with adaptive damped least-squares inversion.
    \item \textbf{Position mode}: The phone's relative displacement from a reference pose is mapped to a target end-effector pose, which is reached via analytical closed-form inverse kinematics.
\end{itemize}

The system operates in simulation, with the robot visualized in RViz~2. Real-time feedback---including manipulability, joint limit proximity, and feasibility---is computed by the kinematics node and sent back to the phone, where it is displayed through visual indicators and haptic feedback.

The project builds directly on the kinematic foundations developed throughout ME/CDS/EE~235a: Denavit-Hartenberg parameterization, forward kinematics, the geometric Jacobian, resolved-rate control, and analytical inverse kinematics for the UR10e. The contribution of this work is the integration of these techniques into a real-time teleoperation pipeline driven by smartphone pose estimation.
